\section*{Capítulo \ref{ch:b3:stable-isotopes} --- Geoquímica de isótopos estáveis}

\begin{solution}{exo:b3:stable-isotopes:1}
$\alpha = 1030.5/997.9 = 1.032669$; $1000\ln\alpha = 32.15$; $\Delta = 30.5 -
(-2.1) = \qty{32.6}{\permil}$. Elas diferem por \qty{0.45}{\permil}, o
$(\delta_A^2 - \delta_B^2)/2000 = 0.46$ da proposição.
\end{solution}

\begin{solution}{exo:b3:stable-isotopes:2}
$\delta^{18}\mathrm O = 0.69\times(-25) - 13.6 = \qty{-30.9}{\permil}$;
$\delta\mathrm D = 5.6\times(-25) - 100 = \qty{-240}{\permil}$; $d = -240 +
8\times30.85 = \qty{6.8}{\permil}$.
\end{solution}

\begin{solution}{exo:b3:stable-isotopes:3}
$\Delta_B = -26.0 - 1.0 = \qty{-27.0}{\permil}$; $f_{\mathrm{org}} = (-5 -
1.0)/(-27.0) = 0.22$.
\end{solution}

\begin{solution}{exo:b3:stable-isotopes:4}
$T^2 = 0.0391\times10^6/(0.600 - 0.154) = \num{87668}$, $T = \qty{296.1}{K}$,
\qty{22.9}{\celsius}.
\end{solution}

\begin{solution}{exo:b3:stable-isotopes:5}
$f^{0.01} = 967/987$, $f = \exp(100\ln0.97974) = 0.129$: \qty{87}{\percent} do
vapor foi removido por chuva. A chuva: $1.0100\times967 - 1000 =
\qty{-23.3}{\permil}$.
\end{solution}

\begin{solution}{exo:b3:stable-isotopes:6}
$\alpha = 1000/960 = 1.041667$; $1000\ln\alpha = 40.82$; $1000(\alpha - 1) =
41.67$; $\Delta = 40$. As duas aproximações erram por cerca de \qty{0.8}{\permil},
a segunda pelo dobro da primeira: com \qty{40}{\permil}, a expansão até a
primeira ordem já não basta.
\end{solution}

\begin{solution}{exo:b3:stable-isotopes:7}
$\delta_c = (1.50 + 29.99)/0.97001 = \qty{32.46}{\permil}$ VSMOW; $1000\ln\alpha
= 1000\ln1.032464 = 31.95$; $T = 18030/(31.95 + 32.42) = \qty{280.1}{K}$,
\qty{7.0}{\celsius}. O fracionamento varia cerca de \qty{0.23}{\permil} por
kelvin, de modo que \qty{0.1}{\permil} na água corresponde a cerca de \qty{0.4}{\celsius}.
\end{solution}

\begin{solution}{exo:b3:stable-isotopes:8}
Volume de gelo $125\times0.011 = \qty{1.375}{\permil}$; o restante, $1.77 - 1.375 =
\qty{0.40}{\permil}$. Perto de \qty{2}{\celsius}, a inclinação é $18030/275.15^2 =
\qty{0.238}{\permil.K^{-1}}$, de modo que a água profunda estava cerca de \qty{1.7}{\celsius}
mais fria.
\end{solution}

\begin{solution}{exo:b3:stable-isotopes:9}
$\Delta^{33}\mathrm S = 8.2 - 0.515\times10.0 = \qty{3.05}{\permil}$, longe de 0:
enxofre com \omterm{def:b3:stable-isotopes:mif}{fracionamento independente da massa}, que só se forma sem uma camada
de ozônio. A pirita formou-se em uma atmosfera sem oxigênio significativo, antes
de cerca de \qty{2.45}{\Ga}.
\end{solution}

\begin{solution}{exo:b3:stable-isotopes:10}
O isótopo pesado se conserva no sistema fechado: $R_0 = fR_v + (1 -
f)\alpha R_v$, logo $R_v/R_0 = 1/[f + \alpha(1 - f)]$. Para $f = 0.2$, $\alpha =
1.01$: $1/1.008$, o vapor fica \qty{7.9}{\permil} mais leve; a \omterm{prop:b3:stable-isotopes:rayleigh-distillation}{destilação de Rayleigh} dá
$0.2^{0.01}$, \qty{16.0}{\permil} mais leve. Remover cada condensado de imediato
empurra o vapor duas vezes mais longe.
\end{solution}

\begin{solution}{exo:b3:stable-isotopes:11}
$M = m(\delta_s - \Delta)/\Delta = 2000\times(-60 + 2.5)/(-2.5) = \qty{46000}{Gt}$.
Então $m = 46000\times(-2.5)/(-25 + 2.5) = \qty{5100}{Gt}$, um pouco abaixo dos
6000 a \qty{12000}{Gt} publicados, cujo limite superior corresponde à excursão
maior.
\end{solution}

\begin{solution}{exo:b3:stable-isotopes:12}
$\delta_v = (1000 + \delta_{v,0})\exp\int(\alpha - 1)\,\mathrm d\ln f - 1000$, logo
$\mathrm d\delta_v = (1000 + \delta_v)(\alpha - 1)\,\mathrm d\ln f$, e $f =
e_s/e_{s,0}$ dá $\mathrm d\ln f = \mathrm d\ln e_s$. A \qty{-20}{\celsius},
$\alpha = 1.01872$: $0.01872\times950.9\times0.0960 =
\qty{1.71}{\permil.K^{-1}}$ para o vapor.
\end{solution}

\begin{solution}{pb:b3:stable-isotopes:1}
\textbf{1.} $\alpha^{18}$: 1.01605, 1.01691, 1.01779, 1.01872, 1.01968,
1.02068; $\alpha^{\mathrm D}$: 1.1424, 1.1524, 1.1630, 1.1744, 1.1866, 1.1997
(a $-5$ a \qty{-30}{\celsius}).

\textbf{2.} $980\times0.425^{0.01605} - 1000 = \qty{-33.4}{\permil}$.

\textbf{3.} $1.01691\times966.6 - 1000 = \qty{-17.0}{\permil}$.

\textbf{4.} Vapor $966.6\times(0.3976)^{0.01779} - 1000 = \qty{-49.1}{\permil}$;
neve $1.01872\times950.9 - 1000 = \qty{-31.3}{\permil}$.

\textbf{5.} Vapor $950.9\times(0.3680)^{0.01968} - 1000 = \qty{-67.6}{\permil}$;
neve $1.02068\times932.4 - 1000 = \qty{-48.3}{\permil}$.

\textbf{6.} $(932.4/980)\times0.0622 = 0.059$: \qty{5.9}{\percent} do
oxigênio-18, contra \qty{6.2}{\percent} da água: o vapor perdeu o isótopo pesado
mais depressa.

\textbf{7.} $(-17.0 + 48.3)/20 = \qty{1.57}{\permil}$ por grau.

\textbf{8.} $t = (-48.3 + 13.6)/0.69 = \qty{-50}{\celsius}$. A relação vale
para temperaturas de superfície; sobre um manto de gelo a superfície é muito mais fria que o
ar onde a neve se forma, e o modelo ignora a mistura e os efeitos cinéticos.

\textbf{9.} Vapor $-247.5$, $-352.5$, \qty{-462.7}{\permil}; neve $\delta\mathrm D
= -132.8$, $-239.6$, \qty{-355.4}{\permil}.

\textbf{10.} $d = -132.8 + 8\times17.02 = 3.3$; depois 10.7 e
\qty{31.3}{\permil}: o excesso cresce para o interior. A inclinação local da
destilação, $\mathrm d\delta\mathrm D/\mathrm d\delta^{18}\mathrm O = (\alpha^{\mathrm D} -
1)(1000 + \delta\mathrm D)/[(\alpha^{18} - 1)(1000 + \delta^{18}\mathrm O)]$, cai abaixo de
8 à medida que o vapor fica muito leve ($1000 + \delta\mathrm D$ diminui bem mais depressa que
$1000 + \delta^{18}\mathrm O$), de modo que a neve se desloca acima da reta.

\textbf{11.} A reta fornece $8\times(-48.34) + 10 = \qty{-376.7}{\permil}$; a
neve fica \qty{21}{\permil} acima dela. O modelo tem apenas fracionamento de equilíbrio;
a neve real se forma em ar supersaturado, com efeito cinético,
e massas de ar se misturam com outras ao longo do caminho.

\textbf{12.} $\delta_c = (3.25 + 29.99)/0.97001 = 34.27$; $1000\ln\alpha =
33.69$; $T = 18030/66.11 = \qty{272.7}{K}$, \qty{-0.4}{\celsius}.

\textbf{13.} $\delta_c = 36.09$, $1000\ln\alpha = 35.46$, $T = \qty{265.6}{K}$,
\qty{-7.5}{\celsius}: mais frio do que pode ficar a água do mar sem congelar. A
variação não pode ser toda devida à temperatura.

\textbf{14.} $\delta_w = 125\times0.011 = \qty{+1.375}{\permil}$; $1000\ln(1036.09/1001.375)
= 34.08$; $T = 18030/66.50 = \qty{271.1}{K}$, \qty{-2.0}{\celsius}: um resfriamento de
\qty{1.6}{\celsius}.

\textbf{15.} $\Delta_{47} = 0.0391\times10^6/272.71^2 + 0.154 = 0.6797$ e
$0.0391\times10^6/271.12^2 + 0.154 = 0.6859$: uma diferença de
\qty{0.006}{\permil}.

\textbf{16.} Ela mede como os isótopos pesados se distribuem dentro do
mineral, um equilíbrio interno que depende apenas da temperatura.

\textbf{17.} $f_{\mathrm{org}} = (-5 - 0.5)/(-28) = 0.20$.

\textbf{18.} $m = 46000\times(-3.0)/(-60 + 3.0) = \qty{2400}{Gt}$.

\textbf{19.} $46000\times(-3.0)/(-25 + 3.0) = \qty{6300}{Gt}$.

\textbf{20.} $46000\times(-3.0)/(-5 + 3.0) = \qty{69000}{Gt}$, mais que o
próprio reservatório, em um intervalo de tempo geologicamente breve: implausível.

\textbf{21.} A profundidade e a extensão da dissolução de carbonato no testemunho, que
mede quanto carbono entrou no oceano, independentemente de sua composição;
junto com o tamanho da excursão, ela fixa $\delta_s$.

\textbf{22.} A neve no Dome Ilse tem $\delta^{18}\mathrm O =
\textbf{\qty[text-series-to-math=true]{-48.3}{\permil}} $ (\qty{-48.4}{\permil}
com $\alpha$ variando continuamente).
\end{solution}
